\documentclass[10pt,a4paper]{amsart}
\usepackage[margin=19mm]{geometry}
\usepackage{amsmath,amssymb,mathtools}
\usepackage[scr=rsfs,cal=euler]{mathalfa}
\usepackage{xfrac}
\usepackage[unicode,hidelinks,bookmarks=false]{hyperref}
\newcommand{\ie}{i.e.\ }
\newcommand{\cf}{cf.\ }
\newcommand{\resp}{respectively\ }
\newcommand{\Subsection}{\subsection}
\providecommand{\nobreakdash}{\nobreak\hbox{-}}
\setlength{\emergencystretch}{2em}
\multlinegap0pt
\usepackage{bm}
\usepackage{bbm}
\usepackage{dsfont}
\usepackage{xspace}
\usepackage{cancel}
\usepackage{mathtools}
\usepackage{amsthm}
\makeatletter
\@ifundefined{lem}{\newtheorem{lem}{Lemma}}{}
\makeatother
\title[Knot exteriors with all compact surfaces of positive genus e]{Knot exteriors with all compact surfaces of positive genus essentially embedded}
\author{Joao M. Nogueira}
\date{Source: arXiv:2508.17504v1 (CC BY 4.0)}
\begin{document}
\maketitle

\noindent\emph{Source and licence.} Knot exteriors with all compact surfaces of positive genus essentially embedded. Version arXiv:2508.17504v1 (CC BY 4.0), distributed under the Creative Commons Attribution 4.0 International licence, as recorded in the arXiv licence metadata for this exact version and shown on the abstract page https://arxiv.org/abs/2508.17504v1. Full rights notice: licenses/arxiv-2508.17504-CC-BY-4.0.txt.\par\smallskip

\noindent\textit{Editorial scope.} Counted unit: the Lem whose source label is \texttt{lemma:incompressible} in the article (source0.tex lines 514--516, source label \texttt{lemma:incompressible}), delivered together with the article's own complete original proof of it (source0.tex lines 517--533, one \texttt{proof} environment of 17 lines). No mathematics is written, repaired or re-derived: statement and proof are the article's own text line for line. Required context retained uncounted: lemma:prime (lines 211--213). Every definition, display and label of those lines is retained unchanged. Omitted from this package and not redistributed: 1--210, 214--513, 534--635. Nothing inside a retained excerpt is omitted or altered: every retained body is delivered exactly as the article's lines are written, so no omission receipt of any kind accompanies this package. Citations: the retained excerpts carry no citation command at all, so no bibliography entry of the article is transcribed and no reference is redistributed. Reference adaptation: none was needed and none was made; every reference inside the delivered unit resolves inside it, so no unresolved reference or citation is produced. Printed slips kept literal: no printed word, symbol, hypothesis or number of the counted statement or of its proof was altered. Layout and class: the article's own class is \texttt{amsart} with packages including amsmath, amsthm, amsfonts, eucal, amssymb, babel, verbatim, inputenc, graphicx, layout, xy, hyperref, caption, subcaption; the wrapper is the lane's pinned \texttt{amsart} editorial wrapper at 10pt on A4, which restates the theorem-like environments the retained text needs and adds the article's own macro definitions that the retained text needs (none), with extra wrapper packages (bm, bbm, dsfont, xspace, cancel, mathtools). The delivered numbering is the unit's own. Assets: no retained span contains a figure environment, a graphics-inclusion command, a PDF-page inclusion, a table or a TikZ picture; no image file is copied into this package, the package declares no asset dependency and no image file is redistributed with it. Licence: version arXiv:2508.17504v1 of this article is distributed under the Creative Commons Attribution 4.0 International licence, as recorded in the arXiv licence metadata for this exact version and shown on the abstract page https://arxiv.org/abs/2508.17504v1. A full rights notice is delivered at licenses/arxiv-2508.17504-CC-BY-4.0.txt. Characters: every retained excerpt is pure ASCII, so the delivered engine, XeLaTeX with its default Latin Modern text font, reports no missing character.
\begin{lem}\label{lemma:prime}
A square double of a non-trivial knot is prime. The only square double of the unknot that is composite is the square knot.
\end{lem}

\begin{lem}\label{lemma:incompressible}
The branched surface $R_g$,  $g\geq 1$,  is incompressible and without Reeb components.
\end{lem}

\begin{proof}
First we observe that $R_g$ doesn't carry a Reeb component.  In fact,  if $R_g$ would carry a torus $T$,  this torus couldn't be transverse to the regular neighborhood of sections of $R_g$ with boundary components, or of sections of $R_g$ of genus bigger than $1$.  In this case,  it could only be transverse to regular neighborhoods of sections $A_1$, $A'$,  $X_g$ (in case $g=1$),  $U_1$,  $V_1$, $U_2$, $V_2$ and $A_2$.  Only one boundary component of $A_1$ and of $A_2$ is connected to these other sections. Hence, $A_1$ and $A_2$ cannot contribute for $R_g$ to carry a torus.  The annulli $U_2$ and $V_2$ are smoothed towards $A_2$.  As $W_{A_2}=0$ in an invariant measure for a torus carried by $R_g$,  we have that necessarily $W_{U_2}=W_{V_2}=0$ as well.  Similarly,  $U_1$, $V_1$ and $X_g$ are smoothed towards $A_1$.  As $W_{A_1}=0$ in an invariant measure for a torus carried by $R_g$,  we also have that  $W_{U_1}=W_{V_1}=W_{X_g}=0$.  Then $A'$ is the only section left whose regular neighborhood could be transverse to a torus.  As $A'$ is smoothed towards $U_1$,  $X_g$ and $V_2$, similarly we have that $W_{A'}=0$ in an invariant measure of a torus carried by $R_g$.  Therefore,  $R_g$ does not carry a torus.\\

If $R_g$ would carry a properly embedded annulus in the exterior of $K$,  $E(K)$,  this annulus would have to be transverse to the regular neighborhood of sections of $R_g$ with boundary components,  that is $P$. However, $P$ has negative Euler characteristics, and the remaining sections of $R_g$ have non-positive Euler characteristics. Hence, any surface carried by $R_g$ with positive weight on $P$ has negative Euler characteristics, which cannot be an annulus.  Therefore, $R_g$ does not carry an annulus.\\

								
Now we prove that $R_g$ is incompressible in $E(K)$. First observe that there are no (half) disks of contact as no circle on the branched locus of $R$ bounds a disk in $\partial_h N(R_g)$ and there are no properly embedded arcs on the branched locus of $R_g$. The space $N(R_g)$ decomposes $E(K)$ into four components: a component cut from $E(K)$ by $X_g$ and $V_1$, denoted $M_1$; a component cut from $E(K)$ by $X_g$ and $U_1$,  denoted $M_1'$; a component cut from $E(K)$ by $U_2$, and $V_2$, denoted $M_2$; a component cut from $E(K)$ by all sections but $X_g$ denoted $M$.\\

There are six components in $\partial_v N(R_g)$: one in $M_1$,  two in $M_1'$, one in $M_2$ and two in $M$.  The component of $\partial_v N(R_g)$ in $M_1$ corresponds to a non-separating circle in $X_g\cup V_1$.  Hence, a monogon in $M_1$  is a compressing disk for $X_g\cup V_1$ in $M_1$,  but this contradicts $X_g$ being essential in the exterior of $J_1$.  A monogon in $M_1'$ is impossible as the boundary of a monogon disk in $M_1'$ would have to intersect one boundary component of the annulus $U_1$ but be disjoint from the other boundary component,  as both correspond to components of $\partial_v N(R_g)$.  Therefore,  there are no monogons in $M_1$ and in $M_1'$.  Similarly, as $\partial_v N(R_g)$ in $M_2$ corresponds to a non-separating circle in $U_2\cup V_2$ and $U_2\cup V_2$ is essential in the exterior of $J_2$,   there are no monogons in $M_2$.\\
Each curve of the branched locus corresponding to the two components of $\partial_v N(R_g)$ in $M$ are non-separating in $\partial M$. Hence,  a monogon in $M$ corresponds to a compressing disk for $\partial M$ in $M$. Let $c_1$ be the curve $V_1\cap A'$ and $c_2$ be the curve $V_2\cap A'$. Suppose there is a monogon $O_1$ corresponding to $c_1$. Then the surface $Y_1$ defined by $P$, $A_1$, $V_1$, $A'$, $V_2$ and $A_2$ has a compressing disk $O_1$.  But $Y_1$ is a genus one surface bounded by $K$, which is incompressible because $K$ is non-trivial and hence we have a contradiction. Suppose there is a monogon $O_2$ corresponding to $c_2$.  Then the surface defined by $P, \, A_1, \, U_1, \, A', \,  U_2$ and $A_2$ has compressing disk $O_2$.  But $Y_2$ is a genus one surface bounded by $K$,  so it is a minimal genus Seifert surface and hence it is essential, contradicting $O_2$ being a compressing disk for $O_2$.  Therefore, there are no monogons in $M$.\\  


We proceed to prove that $\partial_h N(R_g)$ is (boundary) incompressible in the exterior of $K$. In $M_1$,  $\partial_h N(R_g)$ corresponds to the complement of $A_1 \cap X_g \cap V_1$ in $X_g\cup V_1$. Hence,  as $X_g$ is incompressible in $M_1$ we have that $\partial_h N(R_g)\cap M_1$ is incompressible.  In $M_1'$, we argue similarly that  $\partial_h N(R_g)\cap M_1'$ is incompressible.  In $M_2$, $\partial_h N(R_g)\cap M_2$ corresponds to the complement of $A_2 \cap U_2 \cap V_2$ in $U_2\cup V_2$, which corresponds to the exterior in $S^3$ of the knot $J_2$. Hence, $\partial_h N(R_g)\cap M_2$ is incompressible in $M_2$. At last, in $M$, $\partial_h N(R_g)\cap M$ corresponds to the complement of $\partial A'$ in $\partial M$, which is represented by $A'$ and $P$, after an ambient isotopy. An essential simple closed curve in $A'$ corresponds to a meridian of the satellite torus of $K$ (with pattern $J_1\#J_2$), hence there is no compressing disk for $A'$ in $M$. The twice punctured disk $P$ is essential in $M$, otherwise a (boundary) compressing disk would also be a (boundary) compressing disk for the Seifert surface defined by $P\cup A$  (see also the end of proof for Lemma \ref{lemma:prime}).\\


Hence, $\partial_h N(R_g)$ is incompressible and boundary incompressible in $E(K)$, and there are no Reeb components,  monogons or disks of contact  in $R_g$.  Therefore, $R_g$,  $g\geq 1$, is an incompressible branched surface without Reeb components.
\end{proof}

\end{document}
